% Draft of Section 6.1 (replaces the current 6.1, which describes the live market simulator, not the Chapter 6 experiments).
% Settings checked against code and the final matched protocol (docs/chapter6_outline.md, notes/studies/chapter6_metrics.md).
%
% Cross-references: the thesis source was not available, only its text export (Thesis_5). Labels below are placeholders
% named for the sections they must point to; replace them with the thesis's own labels:
%   sec:ch2-tail     -> 2.4.4 Tail Dependence
%   sec:ch4-quant    -> 4.2.1 Quantile Summary Methods
%   sec:ch4-xiao     -> 4.2.2 Counting Matrix Estimator
%   sec:ch4-hermite  -> 4.2.3 Hermite Polynomial Estimator
%   sec:ch4-mom      -> 4.3.1 Rank-Based Method of Moments
%   sec:ch4-mpl      -> 4.3.2 Online Pseudo-Maximum Likelihood Estimation (MPL)
%   sec:ch5-ess      -> 5.1.1 Effective Sample Size
%   sec:ch5-vff      -> 5.2 Variable Forgetting Factor
%   sec:ch5-scoring  -> 5.3.1 Estimator Scoring
%   sec:ch5-bank     -> 5.3.3 Final Bank Algorithm (alg:ch5-bank -> the bank algorithm in 5.3.3)
%   sec:ch5-bounds   -> 5.3.4 Theoretical Bounds
%   app:ch6-design   -> new appendix (docs/chapter6_drafts/appendix_6_design.tex)
% Three methods used here are not yet described in Chapters 4-5 of Thesis_5: the exponentially weighted marginal, the
% common VFF controller and the family bank. This section defines each in one paragraph; the appendix gives details.
\section{Experimental design}\label{sec:ch6-design}

We evaluate the estimators of Chapters 4 and 5 on simulated streams whose copula is known at every time step.
Table~\ref{tab:ch6-experiments} states the five questions and the experiments that answer them.

\begin{table}[ht]
\centering\small
\caption{Questions, comparisons and metrics of the five experiments.}\label{tab:ch6-experiments}
\begin{tabular}{@{}lp{0.40\linewidth}p{0.27\linewidth}p{0.18\linewidth}@{}}
\toprule
 & Question & Comparison & Metrics \\
\midrule
E1 & How do Xiao, Hermite and recursive MPL compare in accuracy and computational cost under limited memory? & Single estimators within 16, 64 and 256\,KiB & $\rho_S$ RMSE, memory, runtime \\
E2 & How does effective memory trade estimation variance against adaptation, and does the approximate optimal memory describe the observed one? & Fixed exponential forgetting over $N_{\mathrm{eff}}$ & $\rho_S$ RMSE \\
E3 & When does adaptive weighting of forgetting factors help? & Fixed memory, no forgetting, VFF, banks & $\rho_S$ RMSE, excess log loss, $H^2$ \\
E4 & How do the scoring loss and the fixed-share step affect the bank? & Squared vs.\ logarithmic loss; Hedge vs.\ fixed share & as E3; weights, regret \\
E5 & Does accurate $\rho_S$ estimation imply accurate copula-density prediction, and do multiple copula families help when the dependence structure changes? & Family bank vs.\ Gaussian forecasts & excess log loss, $H^2$, tail RMSE \\
\bottomrule
\end{tabular}
\end{table}

\subsection{Simulated streams}\label{sec:ch6-dgp}

Each stream contains $T = 5000$ pairs $(X_t, Y_t)$ with standard normal margins and copula $C_t$.
The copula belongs to one of four families: Gaussian, Clayton, Gumbel and Student-$t$ with $\nu = 4$.
A path $d_t$ sets the target Spearman correlation
\begin{equation}\label{eq:ch6-target}
  \rho_{S,t} = \tfrac{6}{\pi}\arcsin(d_t/2),
\end{equation}
the Spearman correlation of a Gaussian copula with correlation $d_t$.
Each family's parameter is the method-of-moments inverse $\theta_t = g_F^{-1}(\rho_{S,t})$ of Section~\ref{sec:ch4-mom}.
The true $\rho_S$ path is therefore the same for all four families, and differences between families reflect copula shape alone.

We use three paths:
\begin{equation}\label{eq:ch6-paths}
  d_t = 0.6 \;\;\text{(constant)}, \qquad
  d_t = 0.2 + 0.6\,\tfrac{t-1}{T-1} \;\;\text{(drift)}, \qquad
  d_t = 0.2 + 0.6\,\mathbb 1\{t > 2500\} \;\;\text{(jump)}.
\end{equation}
These give $\rho_S = 0.582$ on the constant path and $\rho_S$ from $0.191$ to $0.786$ on the other two.

E5 adds four regimes that hold $\rho_S = 0.582$ fixed and change the copula at $t = 2500$:
\begin{itemize}
  \item Gaussian $\to t_4$;
  \item Clayton $\to$ Gumbel;
  \item Gaussian $\to$ Clayton;
  \item a Clayton--Gumbel mixture $\pi_t c^{\mathrm{Cl}} + (1-\pi_t) c^{\mathrm{Gu}}$, with both components at $\rho_S = 0.582$ and Clayton weight $\pi_t$ falling from $0.75$ to $0.25$.
\end{itemize}
Because $\rho_S$ is linear in the copula, the mixture's $\rho_S$ stays fixed while its lower and upper tail dependence move in opposite directions.
No copula available to the family bank equals this mixture.

\subsection{Estimators and comparators}\label{sec:ch6-methods}

\paragraph{Pseudo-observations.}
All banks and comparators share one exponentially weighted empirical marginal per variable, with memory $N_{\mathrm{eff},m} = 500$.
It extends the weighted quantile summary of Section~\ref{sec:ch4-quant} to exponential forgetting (Appendix~\ref{app:ch6-design}).
Before absorbing $X_t$ it returns the pseudo-observation $\widehat U_t$, a forgetting analogue of the rank $(R_t + 1)/(t + 2)$.

\paragraph{Estimators.}
The three estimators are:
\begin{itemize}
  \item Xiao's counting matrix (Section~\ref{sec:ch4-xiao});
  \item the Hermite series estimator (Section~\ref{sec:ch4-hermite});
  \item the recursive Gaussian MPL (Section~\ref{sec:ch4-mpl}).
\end{itemize}
Each forgets exponentially with effective sample size $N_{\mathrm{eff}} \in \{39, 99, 199, 399, 1999\}$ (Section~\ref{sec:ch5-ess}).
Their fixed settings are in Appendix~\ref{app:ch6-design}.

\paragraph{Forgetting-factor banks.}
A bank (Algorithm~\ref{alg:ch5-bank}, Section~\ref{sec:ch5-bank}) combines the five memories of one estimator.
Its estimate is $\widehat\rho_{S,t} = \sum_k w_{k,t}\widehat\rho_{S,k,t}$.
It updates the weights by Hedge with fixed share,
\begin{equation}\label{eq:ch6-hedge}
  w_{k,t} = (1-\delta)\,\frac{w_{k,t-1}e^{-\eta\ell_{k,t}}}{\sum_j w_{j,t-1}e^{-\eta\ell_{j,t}}} + \frac{\delta}{K}, \qquad K = 5.
\end{equation}
The online loss $\ell_{k,t}$ is one of the two scoring rules of Section~\ref{sec:ch5-scoring}, both evaluated at the shared pseudo-observations $(\widehat U_t, \widehat V_t)$:
\begin{itemize}
  \item the squared loss $(p_t - \widehat\rho_{S,k,t-1})^2$ with $p_t = 12(\widehat U_t - \tfrac12)(\widehat V_t - \tfrac12)$;
  \item the logarithmic loss $-\log c^{\mathrm{Ga}}_{\widehat r_{k,t-1}}(\widehat U_t, \widehat V_t)$, where $\widehat r_{k,t-1} = 2\sin(\pi\widehat\rho_{S,k,t-1}/6)$.
\end{itemize}
Both banks use $(\eta, \delta) = (1, 10^{-4})$.
Plain Hedge sets $\delta = 0$, and the uniform control fixes $w_{k,t} = 1/K$.

\paragraph{Comparators.}
\begin{itemize}
  \item \emph{Fixed memory}: the bank's five arms, scored on their own.
  \item \emph{No forgetting}: a cumulative recursive MPL on the same pseudo-observations.
  \item \emph{Variable forgetting} (VFF, Section~\ref{sec:ch5-vff}): a single controller sets the memory of all three estimators.
    It adapts $\log N_{\mathrm{eff}} \in [\log 39, \log 1999]$ by stochastic gradient descent on the one-step squared prediction error of an exponentially weighted mean of $p_t$.
    It runs at rate $0.02$ (slow) or $0.1$ (fast).
\end{itemize}

\paragraph{Family bank (E5).}
The family bank applies \eqref{eq:ch6-hedge} to 24 parametric copulas: eight shapes (Gaussian, $t_4$, Clayton, Gumbel and the rotated Clayton and Gumbel copulas) at memories 99, 399 and 1999.
Each copula is fitted online by Fisher scoring.
The loss is each copula's own log density, with $\eta = 1$ and $\delta = 2.1\times10^{-4}$.
Its forecast density is the weighted mixture of the 24 copulas.

\subsection{Evaluation against the simulation truth}\label{sec:ch6-metrics}

The online losses above train the bank weights.
They use only the estimator's pseudo-observations.
We evaluate every method, whatever its training loss, against the known copula $c_t$ with the same truth-based metrics.
Each metric scores the forecast held before observation $t$, and lower values are better.

\paragraph{Dependence strength.}
Over a window $\mathcal W$ the primary metric is
\begin{equation}\label{eq:ch6-rmse}
  \mathrm{RMSE}_{\rho_S} = \Bigl(|\mathcal W|^{-1}\textstyle\sum_{t\in\mathcal W}(\widehat\rho_{S,t-1} - \rho_{S,t})^2\Bigr)^{1/2}.
\end{equation}

\paragraph{Copula density.}
The primary density metric is the excess logarithmic loss at the true probability integral transform $(U_t, V_t) = (\Phi(X_t), \Phi(Y_t))$:
\begin{equation}\label{eq:ch6-excess}
  \bar\ell^{\,\mathrm{ex}} = |\mathcal W|^{-1}\textstyle\sum_{t\in\mathcal W}\bigl[\log c_t(U_t, V_t) - \log\widehat c_{t\mid t-1}(U_t, V_t)\bigr],
\end{equation}
in nats per observation.
For a normalised forecast $\widehat c_{t\mid t-1}$ its conditional expectation is $D_{\mathrm{KL}}(c_t\,\|\,\widehat c_{t\mid t-1}) \ge 0$, and the logarithmic score is strictly proper.
Both properties need a valid density, so we report \eqref{eq:ch6-excess} only for copula densities.

The secondary metric is the squared Hellinger distance
\begin{equation}\label{eq:ch6-hellinger}
  H_t^2 = 1 - \int_0^1\!\!\int_0^1\sqrt{c_t(u,v)\,\widehat c_{t\mid t-1}(u,v)}\,du\,dv \in [0, 1].
\end{equation}
It is symmetric, bounded and gives less weight to the tails than \eqref{eq:ch6-excess}.
We compute it by quadrature every 50 observations, averaging uniformly in time within each window (Appendix~\ref{app:ch6-design}).

Two forecasts with the same $\rho_S$ can differ greatly in density.
\eqref{eq:ch6-excess} and \eqref{eq:ch6-hellinger} therefore measure something \eqref{eq:ch6-rmse} cannot.

\paragraph{Forecast densities.}
\begin{itemize}
  \item A recursive MPL arm forecasts its fitted Gaussian copula.
  \item A recursive MPL bank yields two forecasts from the same weights:
    \begin{itemize}
      \item the \emph{single Gaussian} at its combined estimate, $c^{\mathrm{Ga}}_{\widehat r_{t-1}}$ with $\widehat r_{t-1} = 2\sin(\pi\widehat\rho_{S,t-1}/6)$;
      \item the \emph{Gaussian mixture} $\sum_k w_{k,t-1}c^{\mathrm{Ga}}_{\widehat r_{k,t-1}}$.
    \end{itemize}
  \item The family bank forecasts its mixture.
\end{itemize}
Xiao's histogram density is normalised but does not have uniform margins, and Hermite's series density is not normalised (Appendix~\ref{app:ch6-design}).
We therefore evaluate Xiao and Hermite by $\rho_S$ RMSE only.

\paragraph{Tails and other measures.}
The tail experiments use the finite-level coefficients of Section~\ref{sec:ch2-tail} at $q = 0.05$:
\begin{equation}\label{eq:ch6-chi}
  \chi_L(q) = C(q,q)/q, \qquad \chi_U(q) = \bigl[1 - 2(1-q) + C(1-q, 1-q)\bigr]/q.
\end{equation}
We report their RMSE against the truth.
The comparators are exponentially weighted corner counters with memories 99, 399 and 1999.

Where relevant we also report:
\begin{itemize}
  \item bank weight paths and the weighted memory $\sum_k w_{k,t}N_{\mathrm{eff},k}$;
  \item regret against the bounds of Section~\ref{sec:ch5-bounds};
  \item retained state in bytes;
  \item update time.
\end{itemize}

\subsection{Comparison protocol}\label{sec:ch6-protocol}

\paragraph{Seeds and tuning.}
Development seeds 1--30 fixed the bank and family-bank constants $(\eta, \delta)$ and the VFF rates.
Evaluation uses seeds 42--101, which played no part in those choices.
E1 tunes each estimator within each memory budget on validation seeds 1000--1039 and tests on seeds 42--141.
We retuned nothing for any evaluation metric.

\paragraph{Windows.}
We discard the first 500 observations and evaluate on $[500, 5000)$.
On jump and switch paths we also report the pre-change window $[500, 2500)$, the post-change window $[2500, 5000)$, and recovery windows $[2500, 2600)$, $[2600, 3000)$ and $[3000, 4000)$.

\paragraph{Paired comparison.}
Within a cell (family and path, and estimator where relevant), every method sees the same 60 streams.
For methods $A$ and $B$ we report the mean and standard error of the per-seed difference $\Delta_s = M_s(A) - M_s(B)$ of a window metric.
$A$ wins if the mean is below $-2$ standard errors, loses if it is above $+2$, and ties otherwise.
We make no multiplicity correction.

The \emph{hindsight-best fixed memory} is the arm with the lowest mean in that cell, chosen separately for each metric.
It is an oracle that no online method can select in advance.

\paragraph{Controlled and tuned comparisons.}
Three comparisons change exactly one component, so their differences are attributable to it:
\begin{itemize}
  \item Hedge vs.\ fixed share ($\delta$ only);
  \item single Gaussian vs.\ Gaussian mixture (forecast construction only);
  \item the Gaussian-only vs.\ full family bank (the set of shapes only).
\end{itemize}
The remaining comparisons set complete procedures against each other, each tuned separately: squared vs.\ logarithmic banks, banks vs.\ VFF, and the E1 estimators.
They identify the better procedure, not the component responsible.
